% GENERATED FILE. Edit canonical lessons, then run npm run generate:books.
% canonical-source-hash: fnv1a64:d8f7066f5abf7b8f
% canonical-lessons: LA-C549-poema, LA-C549-epigramma, LA-C549-versiculus, LA-C549-recitatio, LA-C549-auditor

\chapter{Poems and readings}
\label{ch:la-poems-and-readings}
\hlchaptermodality{549}

\begin{chapteropening}
\textbf{By the end of this chapter:} \emph{I can talk about poems and public readings.}

Complete the last lesson of chapter 549: I can talk about poems and public readings.
\end{chapteropening}

\section[poēma, poēmatis]{poēma, poēmatis — a poem}

\label{lesson:LA-C549-poema}



\begin{quote}
Before the new one: say the Latin for an ornament, then the Latin for ugly.
\end{quote}

\subsection*{You'll want to know: poēma, poēmatis}
\textbf{poēma, poēmatis} — \textquotedblleft{}a poem\textquotedblright{}.

\begin{quote}
Poēta novum poēma amīcīs legit.
\end{quote}

\begin{quote}
\emph{The poet reads his new poem to his friends.}
\end{quote}

Neuter, third declension, from Greek; dative and ablative plural \emph{poēmatīs}. English \emph{poem}.

\begin{cousinweb}
\textbf{poēma} was borrowed from Greek \emph{poiema} \textquotedblleft{}a thing made, a poem\textquotedblright{}, from \emph{poieo} \textquotedblleft{}to make\textquotedblright{} (EtymDB; Lewis \& Short). English \emph{poem} comes from it.
\end{cousinweb}

\begin{grammarlens}[title={What you've built}]
One more word for telling, explaining and describing at B1.
\end{grammarlens}

\subsection*{Your turn}
\begin{itemize}
\raggedright
  \item \emph{Say it:} \emph{poēma, poēmatis}
  \item \emph{Say it:} \emph{poēma, poēmatis}, once more
  \item \emph{Say it:} a sentence of your own with \emph{poēma, poēmatis}
  \item \emph{Recall:} say the Latin for an ornament, then the Latin for ugly, then say \emph{poēma, poēmatis} again
\end{itemize}

\begin{tcolorbox}[breakable,colback=teal!4,colframe=teal!35!black,title={Before you move on}]
What does poēma, poēmatis mean? (\textquotedblleft{}A poem\textquotedblright{}.) Say it once more.
\end{tcolorbox}
\section[epigramma, epigrammatis]{epigramma, epigrammatis — an epigram, a short poem}

\label{lesson:LA-C549-epigramma}



\begin{quote}
Before the new one: say the Latin for ugly, then the Latin for a poem.
\end{quote}

\subsection*{You'll want to know: epigramma, epigrammatis}
\textbf{epigramma, epigrammatis} — \textquotedblleft{}an epigram, a short poem\textquotedblright{}.

\begin{quote}
Hoc epigramma breve et facētum est.
\end{quote}

\begin{quote}
\emph{This epigram is short and witty.}
\end{quote}

Neuter, third declension, from Greek: first an inscription, then a short, pointed poem. English \emph{epigram}.

\begin{cousinweb}
\textbf{epigramma} was borrowed from Greek \emph{epigramma} \textquotedblleft{}inscription\textquotedblright{}, from \emph{epigrapho} \textquotedblleft{}to write upon\textquotedblright{} (EtymDB; Lewis \& Short). English \emph{epigram} comes from it.
\end{cousinweb}

\begin{grammarlens}[title={What you've built}]
One more word for telling, explaining and describing at B1.
\end{grammarlens}

\subsection*{Your turn}
\begin{itemize}
\raggedright
  \item \emph{Say it:} \emph{epigramma, epigrammatis}
  \item \emph{Say it:} \emph{epigramma, epigrammatis}, once more
  \item \emph{Say it:} a sentence of your own with \emph{epigramma, epigrammatis}
  \item \emph{Recall:} say the Latin for ugly, then the Latin for a poem, then say \emph{epigramma, epigrammatis} again
\end{itemize}

\begin{tcolorbox}[breakable,colback=teal!4,colframe=teal!35!black,title={Before you move on}]
What does epigramma, epigrammatis mean? (\textquotedblleft{}An epigram, a short poem\textquotedblright{}.) Say it once more.
\end{tcolorbox}
\section[versiculus, versiculī]{versiculus, versiculī — a little verse, a short line}

\label{lesson:LA-C549-versiculus}



\begin{quote}
Before the new one: say the Latin for a poem, then the Latin for an epigram.
\end{quote}

\subsection*{You'll want to know: versiculus, versiculī}
\textbf{versiculus, versiculī} — \textquotedblleft{}a little verse, a short line\textquotedblright{}.

\begin{quote}
Amīcus mihi paucōs versiculōs scrībit.
\end{quote}

\begin{quote}
\emph{My friend writes me a few little verses.}
\end{quote}

Masculine, a diminutive of \emph{versus} — a line of verse.

\begin{cousinweb}
\textbf{versiculus} is a diminutive of \emph{versus} \textquotedblleft{}a line of writing\textquotedblright{}, from \emph{vertō} \textquotedblleft{}to turn\textquotedblright{} (Lewis \& Short): the turn at the end of a line.
\end{cousinweb}

\begin{grammarlens}[title={What you've built}]
One more word for telling, explaining and describing at B1.
\end{grammarlens}

\subsection*{Your turn}
\begin{itemize}
\raggedright
  \item \emph{Say it:} \emph{versiculus, versiculī}
  \item \emph{Say it:} \emph{versiculus, versiculī}, once more
  \item \emph{Say it:} a sentence of your own with \emph{versiculus, versiculī}
  \item \emph{Recall:} say the Latin for a poem, then the Latin for an epigram, then say \emph{versiculus, versiculī} again
\end{itemize}

\begin{tcolorbox}[breakable,colback=teal!4,colframe=teal!35!black,title={Before you move on}]
What does versiculus, versiculī mean? (\textquotedblleft{}A little verse, a short line\textquotedblright{}.) Say it once more.
\end{tcolorbox}
\section[recitātiō, recitātiōnis]{recitātiō, recitātiōnis — a public reading}

\label{lesson:LA-C549-recitatio}



\begin{quote}
Before the new one: say the Latin for an epigram, then the Latin for a little verse.
\end{quote}

\subsection*{You'll want to know: recitātiō, recitātiōnis}
\textbf{recitātiō, recitātiōnis} — \textquotedblleft{}a public reading\textquotedblright{}.

\begin{quote}
Hodiē poēta recitātiōnem in ātriō facit.
\end{quote}

\begin{quote}
\emph{Today the poet gives a reading in the hall.}
\end{quote}

Feminine, third declension, from \emph{recitō} — to read aloud. Authors read new work to invited listeners. English \emph{recitation}.

\begin{cousinweb}
\textbf{recitātiō} is the action noun of \emph{recitō} \textquotedblleft{}to read aloud\textquotedblright{}, built on \emph{re-} and \emph{citō} \textquotedblleft{}to call\textquotedblright{} (Lewis \& Short). English \emph{recitation} comes from it.
\end{cousinweb}

\begin{grammarlens}[title={What you've built}]
One more word for telling, explaining and describing at B1.
\end{grammarlens}

\subsection*{Your turn}
\begin{itemize}
\raggedright
  \item \emph{Say it:} \emph{recitātiō, recitātiōnis}
  \item \emph{Say it:} \emph{recitātiō, recitātiōnis}, once more
  \item \emph{Say it:} a sentence of your own with \emph{recitātiō, recitātiōnis}
  \item \emph{Recall:} say the Latin for an epigram, then the Latin for a little verse, then say \emph{recitātiō, recitātiōnis} again
\end{itemize}

\begin{tcolorbox}[breakable,colback=teal!4,colframe=teal!35!black,title={Before you move on}]
What does recitātiō, recitātiōnis mean? (\textquotedblleft{}A public reading\textquotedblright{}.) Say it once more.
\end{tcolorbox}
\section[audītor, audītōris]{audītor, audītōris — a listener, hearer}

\label{lesson:LA-C549-auditor}



\begin{quote}
Before the new one: say the Latin for a little verse, then the Latin for a public reading.
\end{quote}

\subsection*{You'll want to know: audītor, audītōris}
\textbf{audītor, audītōris} — \textquotedblleft{}a listener, hearer\textquotedblright{}.

\begin{quote}
Audītōrēs poētae plaudunt.
\end{quote}

\begin{quote}
\emph{The listeners applaud the poet.}
\end{quote}

Masculine, third declension, from \emph{audiō}; also a pupil. English \emph{auditor}.

\begin{cousinweb}
\textbf{audītor} is built on \emph{audiō} \textquotedblleft{}to hear\textquotedblright{} with \emph{-tor} (EtymDB; Lewis \& Short). English \emph{auditor} and \emph{auditorium} come from the family.
\end{cousinweb}

\begin{grammarlens}[title={What you've built}]
One more word for telling, explaining and describing at B1.
\end{grammarlens}

\subsection*{Your turn}
\begin{itemize}
\raggedright
  \item \emph{Say it:} \emph{audītor, audītōris}
  \item \emph{Say it:} \emph{audītor, audītōris}, once more
  \item \emph{Say it:} a sentence of your own with \emph{audītor, audītōris}
  \item \emph{Recall:} say the Latin for a little verse, then the Latin for a public reading, then say \emph{audītor, audītōris} again
\end{itemize}

\begin{tcolorbox}[breakable,colback=teal!4,colframe=teal!35!black,title={Before you move on}]
What does audītor, audītōris mean? (\textquotedblleft{}A listener, hearer\textquotedblright{}.) Say it once more.
\end{tcolorbox}
